% Luc Destombe --- licence CC BY-NC-SA 4.0
% https://creativecommons.org/licenses/by-nc-sa/4.0/deed.fr
% Fichier autonome : compiler avec pdflatex, deux fois (signets, nombre de pages).
\documentclass[a4paper,10pt]{article}
\makeatletter
\newif\iflycee@niveaux
\lycee@niveauxtrue
\newif\iflycee@palatino
\lycee@palatinotrue

\RequirePackage[utf8]{inputenc}
\RequirePackage[T1]{fontenc}
\RequirePackage[french]{babel}
\RequirePackage{etoolbox}

\newcommand{\ShorthandsOff}{\@ifundefined{shorthandoff}{}{\shorthandoff{;:!?}}}
\newcommand{\ShorthandsOn}{\@ifundefined{shorthandon}{}{\shorthandon{;:!?}}}
\ShorthandsOff
\AtBeginDocument{\ShorthandsOn}
\AtBeginEnvironment{tikzpicture}{\ShorthandsOff}

\RequirePackage{geometry}
\RequirePackage{fancyhdr}
\RequirePackage{lastpage}
\RequirePackage{multicol}
\RequirePackage{array}
\RequirePackage{enumitem}
\RequirePackage{xcolor}
\RequirePackage{needspace}

\RequirePackage{amsmath,amssymb}
\RequirePackage{mathpazo}
\DeclareMathSizes{10.95}{10.95}{8}{5.5}
\begingroup\catcode`\_=\active \catcode`\^=\active
\gdef_#1{\sb{\scriptscriptstyle #1}}
\gdef^#1{\sp{\scriptscriptstyle\lycee@moinsepais #1}}
\endgroup
\newcommand\lycee@moins{\mathbin{%
  \setbox\z@\hbox{$\m@th\scriptscriptstyle\lycee@moinsorig$}%
  \dimen@=.55\wd\z@ \advance\dimen@-.5pt
  \hbox to.75\wd\z@{\hss$\m@th\scriptscriptstyle\vcenter{\hbox{%
    \tikz\draw[line width=.5pt,line cap=round](0,0)--(\the\dimen@,0);}}$\hss}}}
\begingroup\lccode`\~=`\- \lowercase{\endgroup\let~}\lycee@moins
\newcommand\lycee@moinsepais{\mathcode`\-="8000 }
\AtBeginDocument{\mathchardef\lycee@moinsorig=\mathcode`\-\relax}
\AtBeginDocument{\catcode`\_=12 \mathcode`\_="8000
  \catcode`\^=12 \mathcode`\^="8000 }
\newcommand\lycee@exposants{%
  \fontdimen13\textfont2=.48\fontdimen6\textfont2
  \fontdimen14\textfont2=.48\fontdimen6\textfont2
  \fontdimen15\textfont2=.48\fontdimen6\textfont2
  \fontdimen13\scriptfont2=.48\fontdimen6\scriptfont2
  \fontdimen14\scriptfont2=.48\fontdimen6\scriptfont2
  \fontdimen15\scriptfont2=.48\fontdimen6\scriptfont2}
\every@math@size\expandafter{\the\every@math@size\lycee@exposants}
\newlength{\retraitcalculs}
\setlength{\retraitcalculs}{2em}

\RequirePackage{tikz}
\usetikzlibrary{arrows.meta,calc,babel,shapes.misc}
\RequirePackage[most]{tcolorbox}
\tcbuselibrary{skins,breakable}

\definecolor{coulDef}{HTML}{1F4E79}
\definecolor{coulProp}{HTML}{7030A0}
\definecolor{coulRem}{HTML}{C55A11}
\definecolor{coulEx}{HTML}{2E7D32}
\definecolor{coulExo}{HTML}{B03A2E}
\definecolor{coulPart}{HTML}{1F4E79}
\definecolor{coulMeth}{HTML}{1B6E4F}
\definecolor{coulCourbe}{HTML}{0B5ED7}
\definecolor{coulCourbeB}{HTML}{C00000}
\definecolor{coulCourbeC}{HTML}{2E7D32}
\definecolor{coulGrille}{HTML}{8C8C8C}
\definecolor{coulRep}{HTML}{1B6E4F}
\definecolor{coulBar}{HTML}{2E7D32}

\newcommand{\R}{\mathbb{R}}
\newcommand{\Rs}{\mathbb{R}^{*}}
\renewcommand{\le}{\leqslant}
\renewcommand{\ge}{\geqslant}
\renewcommand{\approx}{\simeq}
\newcommand{\vect}[1]{\overrightarrow{#1}}
\newcommand{\norme}[1]{\left\lVert #1 \right\rVert}
\newcommand{\intervalle}[2]{\left[#1\,;#2\right]}
\RequirePackage{cancel}
\renewcommand{\CancelColor}{\color{coulExo}}
\newcommand{\fois}[1]{\times{\color{coulExo}#1}}
\newcommand{\barre}[1]{\cancel{#1}}

\tikzset{grille/.style={coulGrille,line width=0.4pt}}
\newcommand{\quadrillage}[4]{\draw[grille,step=1] (#1,#2) grid (#3,#4);}
\tikzset{
  croix/.style={cross out,draw,line width=0.6pt,inner sep=0pt,minimum size=4pt},
  croixdroite/.style={croix,rotate=45,line width=0.8pt},
}
\newcommand{\pt}[3]{\path (#1) node[croix]{} node[#2,font=\small,inner sep=2.5pt]{$#3$};}

\newcommand{\lycee@repere}[4]{%
  \draw[grille,step=1] (#1,#2) grid (#3,#4);
  \draw[-{Stealth[length=2mm]},thick] (#1,0)--({#3+0.4},0) node[below left=-1pt]{$x$};
  \draw[-{Stealth[length=2mm]},thick] (0,#2)--(0,{#4+0.4}) node[below left=-1pt]{$y$};
  \node[below left,font=\scriptsize,inner sep=1.5pt] at (0,0) {$0$};}
\newcommand{\lycee@gradx}[1]{\foreach \k/\l in {#1}{%
  \draw (\k,0.07)--(\k,-0.07) node[below,font=\scriptsize,inner sep=1.5pt]{$\l$};}}
\newcommand{\lycee@grady}[1]{\foreach \k/\l in {#1}{%
  \draw (0.07,\k)--(-0.07,\k) node[left,font=\scriptsize,inner sep=1.5pt]{$\l$};}}
\newcommand{\lycee@intff}[2]{\left[#1\,;#2\right]}
\newcommand{\lycee@intfo}[2]{\left[#1\,;#2\right[}
\newcommand{\lycee@intof}[2]{\left]#1\,;#2\right]}
\newcommand{\lycee@intoo}[2]{\left]#1\,;#2\right[}
\newcommand{\lycee@ens}[1]{\left\{#1\right\}}
\AtBeginDocument{%
  \providecommand{\repere}{\lycee@repere}%
  \providecommand{\gradx}{\lycee@gradx}%
  \providecommand{\grady}{\lycee@grady}%
  \providecommand{\Cf}{\mathcal{C}\sb{\scriptscriptstyle f}}%
  \providecommand{\Cg}{\mathcal{C}\sb{\scriptscriptstyle g}}%
  \providecommand{\intff}{\lycee@intff}%
  \providecommand{\intfo}{\lycee@intfo}%
  \providecommand{\intof}{\lycee@intof}%
  \providecommand{\intoo}{\lycee@intoo}%
  \providecommand{\ens}{\lycee@ens}}

\newlist{questions}{enumerate}{3}
\setlist[questions,1]{label=\arabic*\textdegree),leftmargin=*,itemsep=2pt,topsep=3pt}
\setlist[questions,2]{label=\alph*),leftmargin=*,itemsep=1pt,topsep=2pt}
\setlist[questions,3]{label=\roman*),leftmargin=*,itemsep=1pt,topsep=1pt}
\setlist[itemize]{label=\textbullet,leftmargin=*,itemsep=2pt,topsep=3pt}

\newcommand{\besoinplace}[1]{\par\penalty\@M\begingroup
  \dimen@=#1\relax\dimen@ii=\pagegoal\advance\dimen@ii-\pagetotal
  \ifdim\dimen@>\dimen@ii\ifdim\dimen@ii>\z@\vfil\fi\break\fi\endgroup}

\newcommand{\lycee@pied}{\thepage/\pageref{LastPage}}
\newcommand{\entete}[2]{%
  \pagestyle{fancy}\fancyhf{}%
  \fancyhead[L]{\footnotesize\color{coulPart}\textbf{#1}}%
  \fancyhead[R]{\footnotesize\color{coulPart}\textbf{#2}}%
  \fancyfoot[C]{\footnotesize\lycee@pied}%
  \renewcommand{\headrulewidth}{0.4pt}%
  \renewcommand{\headrule}{\hbox to\headwidth{%
    \color{coulPart!40}\leaders\hrule height \headrulewidth\hfill}}}

\RequirePackage{ccicons}
\newcommand{\lycee@licence}{{\scriptsize\color{gray}%
  \copyright\ Luc Destombe --- \ccbyncsa\ CC BY-NC-SA 4.0}}
\newif\iflycee@licence \lycee@licencetrue
\newcommand{\sanslicence}{\lycee@licencefalse}
\AtBeginDocument{\iflycee@licence\AddToHookNext{shipout/foreground}{%
  \put(0,0){\rlap{\kern\dimexpr1in+\hoffset+\oddsidemargin+\textwidth\relax
    \llap{\raisebox{-\dimexpr1in+\voffset+\topmargin+\headheight+\headsep
      +\textheight+\footskip\relax}{\lycee@licence}}}}}\fi}

\newcommand{\formulesserrees}{%
  \setlength{\abovedisplayskip}{3pt}\setlength{\belowdisplayskip}{3pt}%
  \setlength{\abovedisplayshortskip}{2pt}\setlength{\belowdisplayshortskip}{3pt}}

\setlength{\parindent}{0pt}

\RequirePackage[implicit=false,hidelinks,pdfencoding=auto,psdextra,bookmarksopen,
  bookmarksopenlevel=1,pdfcreator={},pdfproducer={}]{hyperref}
\RequirePackage{bookmark}
\hypersetup{pdfauthor={Luc Destombe},pdfkeywords={CC BY-NC-SA 4.0}}
\newcounter{lycee@signet}
\newcommand{\signet}[3][]{\stepcounter{lycee@signet}%
  \edef\lycee@dest{\if\relax\detokenize{#1}\relax
    signet.\arabic{lycee@signet}\else#1\fi}%
  \hypertarget{\lycee@dest}{}\bookmark[level=#2,dest=\lycee@dest]{#3}}
\pdfstringdefDisableCommands{%
  \def\R{R}\def\vect#1{#1}\def\Cf{Cf}\def\Cg{Cg}\def\fois#1{×#1}%
  \def\barre#1{#1}\def\intervalle#1#2{[#1;#2]}\def\intff#1#2{[#1;#2]}%
  \def\textsuperscript#1{#1}\def\dfrac#1#2{#1/#2}\def\frac#1#2{#1/#2}%
  \def\vec#1{#1⃗}}%
\RequirePackage[official]{eurosym}

\geometry{top=0.8cm,bottom=1.3cm,left=1cm,right=1cm,footskip=0.6cm}

\pagestyle{fancy}
\fancyhf{}
\renewcommand{\headrulewidth}{0pt}
\fancyfoot[C]{\footnotesize Page \thepage{} / \pageref{LastPage}}

\newcommand{\titrefiche}[2]{%
  \begin{center}
    \Large\textbf{#1}\\[2pt]
    \large\textbf{#2}\\[2pt]
  \end{center}
  \vspace{0.10cm}}

\setlength{\columnseprule}{0.4pt}
\setlength{\columnsep}{15pt}
\renewcommand{\columnseprulecolor}{\color{black!45}}
\setlength{\multicolsep}{2pt}

\newcounter{partiefiche}
\renewcommand{\thepartiefiche}{\Roman{partiefiche}}
\newif\ifapresbandeau
\newcommand{\lycee@bandeau}[1]{%
  \par\penalty-600
  \vspace{0.08cm}%
  \noindent\colorbox{coulPart}{%
    \parbox{\dimexpr\linewidth-2\fboxsep\relax}{%
      \color{white}\bfseries\raggedright\strut#1\strut}}%
  \nopagebreak\par\nopagebreak
  \vspace{0.06cm}\nopagebreak
  \apresbandeautrue}
\newcommand{\partiefiche}[1]{%
  \refstepcounter{partiefiche}%
  \lycee@bandeau{\signet{1}{\thepartiefiche{} --- #1}\thepartiefiche{} --- #1}}
\newcommand{\partiefichesansnum}[1]{\lycee@bandeau{\signet{1}{#1}#1}}

\newcounter{exo}
\newlength{\espaceexo}\setlength{\espaceexo}{7pt}
\newlength{\espacetitre}\setlength{\espacetitre}{2pt}
\newcommand{\exo}[1][\relax]{%
  \par
  \ifapresbandeau\apresbandeaufalse\else\penalty-150\addvspace{\espaceexo}\fi
  \refstepcounter{exo}%
  \noindent\signet[exercice-\arabic{exo}]{2}{Exercice \arabic{exo}}%
  \textbf{\color{coulExo}\underline{Exercice \theexo}}%
  \ifx\relax#1\relax\else\textbf{ : #1}\fi
  \par\nopagebreak\vspace{\espacetitre}\nopagebreak
  \ignorespaces}

\setlist[enumerate]{nosep,leftmargin=*,itemsep=2pt,topsep=2pt}
\setlist[itemize]{nosep,leftmargin=*,topsep=2pt}
\newcommand{\choix}[3]{\par\hspace*{1em}%
  \makebox[0.3\linewidth][l]{a) #1}\makebox[0.3\linewidth][l]{b) #2}c) #3}
\newcommand{\deux}[2]{\par\hspace*{0.6em}%
  \makebox[0.48\linewidth][l]{#1}#2}
\newcommand{\trois}[3]{\par\hspace*{0.6em}%
  \makebox[0.32\linewidth][l]{#1}\makebox[0.32\linewidth][l]{#2}#3}

  \renewcommand{\exo}[1][\relax]{%
    \ifapresbandeau\apresbandeaufalse\else\par\penalty-150\fi
    \refstepcounter{exo}%
    \vspace{0.05cm}%
    \noindent\signet[exercice-\arabic{exo}]{2}{Exercice \arabic{exo}}%
    \textbf{\color{coulExo}\underline{Exercice \theexo}}%
    \ifx\relax#1\relax\else\textbf{ : #1}\fi%
    \par\nopagebreak
    \ignorespaces}
  \newcommand{\rappel}[1]{%
    \par\vspace{0.06cm}\nopagebreak
    \noindent\fcolorbox{black!35}{black!5}{%
      \parbox{\dimexpr\linewidth-2\fboxsep-2\fboxrule\relax}{%
        \small\textbf{Rappel.} #1}}%
    \par\vspace{0.06cm}\nopagebreak}
  \setlist[enumerate]{nosep,leftmargin=*}
  \setlist[itemize]{nosep,leftmargin=*}
  \setlength{\multicolsep}{3pt plus 1pt minus 1pt}
  \setlength{\premulticols}{10pt}
  \setlength{\postmulticols}{10pt}
  \newlist{expr}{enumerate}{1}
  \setlist[expr]{label={\Alph*\ $=$},nosep,leftmargin=*,itemsep=0pt}

\renewenvironment{center}{\par\addvspace{3pt}\centering}{\par\addvspace{3pt}}
\formulesserrees
\AtBeginDocument{\formulesserrees}
\clubpenalty=10000
\widowpenalty=10000
\displaywidowpenalty=10000
\brokenpenalty=10000
\setlength{\parskip}{0pt}

\RequirePackage{ragged2e}
\setlength{\RaggedRightRightskip}{0pt plus 1fil}
\AtBeginDocument{\RaggedRight\binoppenalty=10000 \relpenalty=10000 }
\g@addto@macro\@parboxrestore{\RaggedRight}
\makeatother
\geometry{top=0.8cm,bottom=1.3cm,left=1cm,right=1cm,footskip=0.7cm,headheight=12pt}

\usepackage{pgfplots}
\pgfplotsset{compat=1.18}
\definecolor{coulCourbeD}{HTML}{7030A0}
\definecolor{coulCourbeE}{HTML}{B8860B}
\pgfplotsset{
  repere/.style={
    axis lines=middle,
    axis line style={-{Stealth[length=2mm]},thick},
    label style={font=\scriptsize},
    tick label style={font=\scriptsize},
    grid=both,
    major grid style={grille},
    minor tick num=0,
    tick align=inside,
    clip mode=individual,
  },
}

\begin{document}
\titrefiche{1\textsuperscript{re} spécialité --- Le second degré}{Fiche d'exercices du chapitre}

\begin{multicols}{2}

\partiefiche{Forme canonique}

\exo[Compléter un carré ($a=1$)]
Écrire chaque trinôme sous la forme $(x-\alpha)^{2}+\beta$.
\begin{multicols}{2}
\begin{expr}
  \item $x^{2}+8x-3$
  \item $x^{2}-6x+2$
  \item $x^{2}-3x+1$
  \item $x^{2}+5x+2$
  \item $x^{2}-12x+35$
  \item $x^{2}-9x$
\end{expr}
\end{multicols}

\exo[Compléter un carré ($a\ne 1$)]
Écrire chaque trinôme sous la forme $a(x-\alpha)^{2}+\beta$.
\begin{multicols}{2}
\begin{expr}
  \item $2x^{2}+8x+3$
  \item $-x^{2}+3x+2$
  \item $4x^{2}-8x+4$
  \item $-x^{2}+5x-4$
  \item $2x^{2}+2x-1$
  \item $-2x^{2}+8x-5$
\end{expr}
\end{multicols}

\exo[Forme canonique par $\alpha$ et $\beta$]
Pour chaque trinôme, calculer $\alpha=-\dfrac{b}{2a}$ et $\beta$, écrire la forme
canonique, puis la développer pour vérifier.
\begin{enumerate}[label=\arabic*)]
  \item $P(x)=x^{2}-6x+1$
  \item $Q(x)=3x^{2}+6x-2$
  \item $R(x)=-x^{2}-4x-9$
\end{enumerate}

\exo[De la forme canonique à la factorisation]
Écrire chaque trinôme sous forme canonique. Le factoriser avec $u^{2}-v^{2}=(u-v)(u+v)$
quand c'est possible ; sinon, expliquer pourquoi il ne se factorise pas.
\begin{multicols}{2}
\begin{expr}
  \item $x^{2}-4x-5$
  \item $x^{2}+4x+7$
  \item $2x^{2}+x-3$
  \item $-x^{2}-2x-3$
\end{expr}
\end{multicols}

\partiefiche{Équations du second degré}

\exo[Factoriser plutôt que calculer $\Delta$]
Résoudre chaque équation en factorisant, sans discriminant.
\begin{multicols}{2}
\begin{enumerate}[label=\arabic*)]
  \item $9x^{2}-4=0$
  \item $x^{2}-6x=0$
  \item $2x^{2}-18=0$
  \item $x^{2}-5=0$
  \item $(3x-2)(x+5)=0$
  \item $x^{2}-10x+25=0$
\end{enumerate}
\end{multicols}

\exo[Avec le discriminant]
Résoudre dans $\R$ en calculant le discriminant.
\begin{enumerate}[label=\arabic*)]
  \item $x^{2}+x-12=0$
  \item $x^{2}+3x+5=0$
  \item $-x^{2}+6x-9=0$
  \item $t^{2}-4t-21=0$
  \item $2-y-3y^{2}=0$
  \item $-2x^{2}+5x+1=0$
\end{enumerate}

\exo[Des racines carrées dans les coefficients]
Résoudre dans $\R$ :
\begin{enumerate}[label=\arabic*)]
  \item $x^{2}-2\sqrt{5}\,x-4=0$
  \item $\sqrt{3}\,t^{2}-4t+\sqrt{3}=0$
  \item $x^{2}-(1+\sqrt{2})x+\sqrt{2}=0$
  \item $4x-x^{2}-5=0$
  \item $x^{3}+2x^{2}-15x=0$
  \item $(3x+2)^{2}+1=0$
\end{enumerate}

\exo[Une parabole et une droite]
\begin{enumerate}[label=\arabic*)]
  \item Dans un même repère, tracer l'allure de la parabole d'équation $y=x^{2}$ et la
    droite d'équation $y=2x+3$.
  \item Expliquer pourquoi les solutions de $x^{2}-2x-3=0$ sont les abscisses des points
    communs aux deux courbes. Combien cette équation a-t-elle de solutions ? Les lire.
  \item Résoudre $x^{2}-2x-3=0$ par le calcul.
\end{enumerate}

\exo[Un programme Python]
La fonction \texttt{racines} renvoie la liste des solutions de $ax^{2}+bx+c=0$.

\footnotesize
\begin{verbatim}
from math import sqrt

def racines(a, b, c):
    d = b*b - 4*a*c
    if d < 0:
        return []
    if d == 0:
        return [-b/(2*a)]
    r = sqrt(d)
    return [(-b - r)/(2*a), (-b + r)/(2*a)]
\end{verbatim}
\normalsize
\begin{enumerate}[label=\arabic*)]
  \item Que représente la variable \texttt{d} ? Que renvoie la fonction selon son signe ?
  \item Saisir ce programme, puis le tester sur les six équations de l'exercice 6 :
    retrouve-t-on les résultats obtenus à la main ?
  \item Que se passe-t-il avec \texttt{racines(0, 3, -12)} ? Pourquoi ? Quelle est
    pourtant la solution de $3x-12=0$ ?
  \item Comparer \texttt{racines(1, 6, 9)} et \texttt{racines(0.2, 1.2, 1.8)}. Le second
    trinôme vaut $0{,}2(x+3)^{2}$ : quelle réponse attendait-on ? Pourquoi un test
    d'égalité est-il fragile avec des nombres à virgule ?
\end{enumerate}

\partiefiche{Factorisation du trinôme}

\exo[Factoriser un trinôme]
Écrire chaque trinôme comme un produit de facteurs.
\begin{enumerate}[label=\arabic*)]
  \item $f(x)=x^{2}+3x-10$
  \item $g(x)=3x^{2}-10x+3$
  \item $h(x)=-2x^{2}+3x+5$
  \item $k(x)=-\dfrac{1}{3}x^{2}+\dfrac{1}{3}x+2$
\end{enumerate}

\exo[Factoriser, ou prouver que c'est impossible]
Pour chaque trinôme, donner une forme factorisée ou démontrer qu'il ne se factorise pas.
\begin{multicols}{2}
\begin{expr}
  \item $x^{2}+2x+2$
  \item $x^{2}+8x+16$
  \item $3x^{2}-5x+2$
  \item $-x^{2}+4x-7$
  \item $9x^{2}+12x+4$
  \item $2x^{2}+3x+1$
\end{expr}
\end{multicols}

\exo[Simplifier un quotient]
Donner l'ensemble de définition, factoriser le numérateur et le dénominateur, puis
simplifier.
\begin{enumerate}[label=\arabic*),itemsep=7pt]
  \item $A(x)=\dfrac{x^{2}+x-6}{x+3}$
  \item $B(x)=\dfrac{3x^{2}+x-2}{x^{2}-1}$
\end{enumerate}

\partiefiche{Somme et produit des racines}

\exo[Une racine connue (1)]
\begin{enumerate}[label=\arabic*)]
  \item Montrer que $-2$ est solution de $x^{2}+5x+6=0$.
  \item Donner la somme $S$ et le produit $P$ des racines.
  \item En déduire l'autre solution, et la contrôler avec $P$.
\end{enumerate}

\exo[Une racine connue (2)]
\begin{enumerate}[label=\arabic*)]
  \item Montrer que $3$ est solution de $x^{2}-7x+12=0$.
  \item Donner la somme $S$ et le produit $P$ des racines.
  \item En déduire l'autre solution, et la contrôler avec $P$.
\end{enumerate}

\exo[Racine évidente, puis somme]
Pour chaque équation, trouver une racine évidente parmi $1$, $-1$, $2$, $-2$\dots{} puis
l'autre racine avec la somme $S=-\dfrac{b}{a}$.
\begin{multicols}{2}
\begin{enumerate}[label=\arabic*)]
  \item $x^{2}-6x+5=0$
  \item $4x^{2}+x-5=0$
  \item $7x^{2}-9x+2=0$
  \item $3x^{2}+8x+5=0$
\end{enumerate}
\end{multicols}

\exo[Sans discriminant]
Résoudre chaque équation sans calculer de discriminant : racine évidente, puis somme ou
produit.
\begin{multicols}{2}
\begin{enumerate}[label=\arabic*)]
  \item $x^{2}-10x+9=0$
  \item $-2x^{2}+3x+5=0$
  \item $x^{2}+4x-12=0$
  \item $x^{2}+7x+6=0$
  \item $x^{2}-9x+18=0$
  \item $x^{2}+x\sqrt{2}-4=0$
\end{enumerate}
\end{multicols}

\exo[Calculer sans résoudre]
On note $x_{1}$ et $x_{2}$ les racines de $2x^{2}-9x+4=0$, sans les calculer.
\begin{enumerate}[label=\arabic*)]
  \item Justifier que ces racines existent, puis donner $S=x_{1}+x_{2}$ et $P=x_{1}x_{2}$.
  \item Calculer $\dfrac{1}{x_{1}}+\dfrac{1}{x_{2}}$ à l'aide de $S$ et $P$.
  \item Développer $S^{2}$, puis en déduire $x_{1}^{2}+x_{2}^{2}$.
\end{enumerate}

\exo[Signe des racines]
Sans calculer les racines, donner leur signe à l'aide de $S$ et $P$.
\begin{enumerate}[label=\arabic*)]
  \item $x^{2}-9x+14=0$
  \item $x^{2}+2x-15=0$
  \item $3x^{2}+10x+3=0$
\end{enumerate}

\partiefiche{Signe du trinôme et inéquations}

\exo[Tableau de signes]
Dresser le tableau de signes de chaque trinôme.
\begin{multicols}{2}
\begin{expr}
  \item $x^{2}-6x+8$
  \item $-3x^{2}+9x-6$
  \item $x^{2}-4x+6$
  \item $-x^{2}-4x-4$
\end{expr}
\end{multicols}

\exo[Inéquations du second degré]
Résoudre :
\begin{multicols}{2}
\begin{enumerate}[label=\arabic*)]
  \item $x^{2}-7x+10>0$
  \item $x^{2}+2x-24<0$
  \item $p^{2}-p-30\le 0$
  \item $x^{2}+3x+4<0$
  \item $2x^{2}-12x+18>0$
  \item $-3x^{2}+2x+2>0$
\end{enumerate}
\end{multicols}

\exo[Position d'une parabole et d'une droite]
On note $\mathcal{C}_{f}$ la courbe de $f(x)=x^{2}+2x+3$ et $\mathcal{D}_{g}$ la droite
d'équation $y=4x+2$.
\begin{enumerate}[label=\arabic*)]
  \item Résoudre $f(x)=4x+2$. Que peut-on dire des points communs à $\mathcal{C}_{f}$ et
    $\mathcal{D}_{g}$ ?
  \item Étudier le signe de $f(x)-(4x+2)$.
  \item En déduire la position de $\mathcal{C}_{f}$ par rapport à $\mathcal{D}_{g}$.
\end{enumerate}

\partiefiche{Représentation graphique}

\exo[Sommet, racines, allure]
Pour chaque fonction : orientation de la parabole, sommet $S(\alpha\,;\beta)$, racines
éventuelles, puis allure de la courbe dans un repère.
\begin{multicols}{2}
\begin{expr}
  \item $x^{2}+4x+3$
  \item $-x^{2}+4x-3$
  \item $x^{2}-4x+5$
  \item $-2x^{2}-4x$
\end{expr}
\end{multicols}

\exo[Utiliser la symétrie]
Soit $f$ définie sur $\R$ par $f(x)=x^{2}-6x+8$.
\begin{enumerate}[label=\arabic*)]
  \item Calculer $\alpha$, $\beta$ et les racines de $f$.
  \item Calculer $f(0)$. Quel point de la parabole est le symétrique de $(0\,;f(0))$ par
    rapport à la droite $x=\alpha$ ? Le vérifier par un calcul.
  \item Placer ces points dans un repère, puis tracer l'allure de la parabole.
\end{enumerate}

\exo[Retrouver un trinôme]
Une parabole coupe l'axe des abscisses en $-1$ et en $3$ ; son sommet est $S(1\,;-8)$.
\begin{enumerate}[label=\arabic*)]
  \item Vérifier que l'abscisse du sommet est le milieu des racines.
  \item Le trinôme s'écrit $a(x+1)(x-3)$. Trouver $a$ à l'aide du sommet.
  \item Écrire le trinôme sous forme développée, et contrôler avec le sommet.
\end{enumerate}

\exo[Le jet d'une fontaine]
Le jet d'eau d'une fontaine suit une parabole. Dans le repère ci-dessous, $x$ est la
distance horizontale (en m) et $y$ la hauteur du jet au-dessus du bassin (en m). Cette
parabole est la courbe d'un trinôme $P(x)=ax^{2}+bx+c$.

\begin{center}
\begin{tikzpicture}
\begin{axis}[repere,
    x=0.7cm,y=0.7cm,
    xmin=-0.6,xmax=6.6,ymin=-0.6,ymax=2.9,
    xtick={1,...,6},ytick={1,2},
    xlabel={$x$},ylabel={$y$}]
  \draw[coulCourbe,line width=1.1pt] (axis cs:1,0)
    .. controls (axis cs:1.667,1.333) and (axis cs:2.333,2) .. (axis cs:3,2)
    .. controls (axis cs:3.667,2) and (axis cs:4.333,1.333) .. (axis cs:5,0);
\end{axis}
\end{tikzpicture}
\end{center}
\begin{enumerate}[label=\arabic*)]
  \item Lire les racines $x_{1}$ et $x_{2}$ de $P$. Que représentent-elles pour le jet ?
  \item Lire les coordonnées $(\alpha\,;\beta)$ du sommet. Que représente $\beta$ ?
  \item Vérifier que $\alpha$ est le milieu de $x_{1}$ et $x_{2}$.
  \item Quel est le signe de $a$ ? Justifier.
\end{enumerate}

\exo[À chaque parabole sa fonction]
Le graphique ci-dessous montre $5$ paraboles $\mathcal{C}_{1}$, \dots, $\mathcal{C}_{5}$.
Associer à chaque fonction sa parabole, en justifiant.
\begin{enumerate}[label=\alph*)]
  \item $f_{1}(x)=-x^{2}-2x-2$
  \item $f_{2}(x)=x^{2}-2x+2$
  \item $f_{3}(x)=2x^{2}+5x+2$
  \item $f_{4}(x)=-2x^{2}+3x+2$
  \item $f_{5}(x)=x^{2}-x+\dfrac{1}{4}$
\end{enumerate}

\begin{center}
\begin{tikzpicture}
\begin{axis}[repere,
    x=0.55cm,y=0.55cm,
    xmin=-5,xmax=5,ymin=-5,ymax=7,
    xtick={-5,...,5},ytick={-5,...,7},
    legend style={font=\scriptsize,at={(0.5,-0.18)},anchor=north,legend columns=5}]
  \addplot[coulCourbe,line width=1pt,domain=-1.25:2.75,samples=100]{-2*x^2+3*x+2};
  \addlegendentry{$\mathcal{C}_1$}
  \addplot[coulCourbeB,line width=1pt,domain=-1.4:3.4,samples=100]{x^2-2*x+2};
  \addlegendentry{$\mathcal{C}_2$}
  \addplot[coulCourbeC,line width=1pt,domain=-3.05:1.05,samples=100]{-x^2-2*x-2};
  \addlegendentry{$\mathcal{C}_3$}
  \addplot[coulCourbeD,line width=1pt,domain=-2.1:3.1,samples=100]{x^2-x+0.25};
  \addlegendentry{$\mathcal{C}_4$}
  \addplot[coulCourbeE,line width=1pt,domain=-3.25:0.75,samples=100]{2*x^2+5*x+2};
  \addlegendentry{$\mathcal{C}_5$}
\end{axis}
\end{tikzpicture}
\end{center}

\partiefiche{Variations de la fonction trinôme}

\exo[Variations et forme canonique]
Dresser le tableau de variations de chaque fonction, définie sur $\R$.
\begin{enumerate}[label=\arabic*)]
  \item $f(x)=3(x-2)^{2}+1$
  \item $g(x)=-2(x+3)^{2}+4$
  \item $h(x)=x(x+6)$
\end{enumerate}

\exo[Variations par $\alpha$ et $\beta$]
Pour chaque trinôme, calculer $\alpha$ et $\beta$, puis dresser son tableau de
variations.
\begin{multicols}{2}
\begin{expr}
  \item $x^{2}+6x+7$
  \item $-x^{2}+4x+5$
  \item $3x^{2}-12x+5$
  \item $-\dfrac{1}{2}x^{2}-2x$
\end{expr}
\end{multicols}

\exo[Un maximum]
Soit $f$ définie sur $\R$ par $f(x)=-3x^{2}+12x-10$.
\begin{enumerate}[label=\arabic*)]
  \item Écrire $f(x)$ sous forme canonique.
  \item En déduire que $f(x)\le 2$ pour tout réel $x$. Pour quelle valeur de $x$
    a-t-on $f(x)=2$ ?
\end{enumerate}

\exo[Étude complète d'un trinôme]
Soit $f$ définie sur $\R$ par $f(x)=-x^{2}+2x+3$.
\begin{enumerate}[label=\arabic*)]
  \item Trouver les racines de $f$, puis écrire $f(x)$ sous forme factorisée.
  \item Calculer $\alpha$ et $\beta$, puis dresser le tableau de variations de $f$.
  \item Vérifier que $\alpha$ est le milieu des deux racines.
  \item Dresser le tableau de signes de $f$, puis tracer l'allure de sa courbe.
\end{enumerate}

\partiefiche{Équations se ramenant au second degré}

\exo[Développer, puis résoudre]
Résoudre dans $\R$ :
\begin{multicols}{2}
\begin{enumerate}[label=\arabic*)]
  \item $(x+1)(x-4)=6$
  \item $(3x+1)^{2}=16$
  \item $x(x-4)=-3$
  \item $(2x+3)^{2}=(x+3)^{2}$
  \item $(x-3)^{2}=x+3$
  \item $(x+4)(x-4)=6x+11$
\end{enumerate}
\end{multicols}

\exo[Un seul quotient]
Donner l'ensemble de définition, puis résoudre :
\begin{multicols}{2}
\begin{enumerate}[label=\arabic*),itemsep=7pt]
  \item $\dfrac{12}{x}=x+1$
  \item $\dfrac{6}{x-1}=x$
  \item $\dfrac{x+10}{x}=x-2$
  \item $x+\dfrac{8}{x}=6$
\end{enumerate}
\end{multicols}

\exo[Équations avec quotients]
Résoudre dans $\R$ :
\begin{enumerate}[label=\arabic*),itemsep=7pt]
  \item $\dfrac{x^{2}-6x+9}{x-3}=2x-5$
  \item $\dfrac{2}{x}+\dfrac{2}{x-3}=1$
  \item $\dfrac{1}{x+2}-\dfrac{2}{x-3}=\dfrac{3}{2}$
  \item $\dfrac{2x^{2}+9x+10}{x+2}=x+3$
\end{enumerate}

\partiefiche{Inéquations se ramenant au second degré}

\exo[Développer, puis résoudre]
Résoudre :
\begin{multicols}{2}
\begin{enumerate}[label=\arabic*)]
  \item $(x+1)(x+3)\le 2x+6$
  \item $(3x+1)^{2}<(x+3)^{2}$
  \item $x(x-4)\ge 5$
  \item $(x+2)^{2}>9$
\end{enumerate}
\end{multicols}

\exo[Signe d'un quotient]
Donner l'ensemble de définition, puis résoudre :
\begin{enumerate}[label=\arabic*),itemsep=7pt]
  \item $\dfrac{2x-6}{x+1}\le 0$
  \item $\dfrac{x^{2}-1}{x+2}\ge 0$
  \item $\dfrac{x^{2}+x-6}{x^{2}-4x+3}>0$
  \item $\dfrac{x^{2}-5x+4}{x-2}<0$
\end{enumerate}

\exo[Tout dans un membre]
Résoudre $\dfrac{x+2}{2-x}\ge 3x+1$.

\emph{On ne multipliera pas par un dénominateur dont on ignore le signe.}

\partiefiche{Problèmes}

\emph{Pour chaque problème : choisir l'inconnue et dire ce qu'elle représente, mettre en
équation, résoudre, puis revenir à l'énoncé et répondre par une phrase.}

\exo[Le nombre mystère]
Le carré d'un nombre positif, diminué de son triple, vaut $40$. Quel est ce nombre ?

\exo[Une voile d'ombrage]
Une voile d'ombrage a la forme d'un triangle rectangle. L'un des côtés de l'angle droit
mesure $2$ m de plus que l'autre, et l'aire de la voile est de $12\text{ m}^{2}$. Quelles
sont les longueurs des côtés de l'angle droit ?

\exo[Agrandir une photo]
On agrandit une photo carrée en ajoutant $12$ cm à son côté : son aire est alors
multipliée par $4$. Quel était le côté de la photo de départ ?

\exo[La bordure d'une tablette]
L'écran d'une tablette mesure $24$ cm sur $16$ cm. Il est entouré d'une bordure de largeur
constante $x$. La face avant de la tablette a une aire de $560\text{ cm}^{2}$.

\begin{center}
\begin{tikzpicture}[scale=0.8,every node/.style={font=\scriptsize}]
  \draw[fill=black!8,rounded corners=2pt] (0,0) rectangle (4.4,3.2);
  \draw[fill=white] (0.4,0.4) rectangle (4,2.8);
  \node at (2.2,1.6) {écran $24$ cm $\times\ 16$ cm};
  \draw[<->] (0,-0.25) -- (0.4,-0.25) node[midway,below]{$x$};
  \draw[<->] (-0.25,2.8) -- (-0.25,3.2) node[midway,left]{$x$};
\end{tikzpicture}
\end{center}

Trouver la largeur $x$ de la bordure.

\exo[Des vélos reconditionnés]
Un atelier remet en état et vend $x$ vélos électriques par mois, avec $0\le x\le 40$. Son
bénéfice mensuel, en euros, est
\[
  B(x)=-10x^{2}+400x-3000
\]
\begin{enumerate}[label=\arabic*)]
  \item Calculer $B(5)$ et $B(12)$, et interpréter.
  \item Résoudre $B(x)>0$. Combien de vélos l'atelier doit-il vendre pour être
    bénéficiaire ?
  \item Pour combien de vélos vendus le bénéfice est-il maximal ? Que vaut-il ?
\end{enumerate}

\vfill

\end{multicols}
\end{document}
